\documentclass[11pt,a4paper]{article}
\usepackage[utf8]{inputenc}
\usepackage[spanish]{babel}
\usepackage[a4paper,left=2.5cm,right=2.5cm,top=2.5cm,bottom=2.5cm]{geometry}
\usepackage{amsmath}
\usepackage{amsfonts}
\usepackage{amssymb}
\usepackage{fancyhdr}
\usepackage{graphicx}
\usepackage{float}
\usepackage{booktabs}
\usepackage{array}
\usepackage{xcolor}
\usepackage{hyperref}
\usepackage{enumitem}
\usepackage{tcolorbox}

\pagestyle{fancy}
\lhead{\small Econometría - Pauta Ejercicios 7}
\rhead{\small Universidad Técnica Federico Santa María}

\definecolor{verdeUSM}{RGB}{0,100,50}


\begin{document}

\begin{center}
    {\Large \textbf{Pauta Ejercicios 5: Econometría}}\\[0.3cm]
    Profesor: Pedro Comte Hidalgo \\
    Ayudante: Jociney Honorio \\
    \textit{Tests F, test de grupos y modelos con dummies e interacciones}
\end{center}

\vspace{0.3cm}
\hrule
\vspace{0.5cm}

\section{Ejercicios}
%=====================================================
\begin{enumerate}
    \item 
Un investigador estima con $n=120$ observaciones el modelo:
\[
y = \beta_0 + \beta_1 x_1 + \beta_2 x_2 + \beta_3 x_3 + \beta_4 x_4 + \beta_5 x_5 + u
\]
obteniendo $R^2 = 0.78$ y $\bar{R}^2 = 0.7706$.

Sospecha que $x_4$ y $x_5$ no aportan. Al estimar el modelo restringido (sin $x_4$, $x_5$) obtiene $R^2_R = 0.74$. Conteste con significancia del 5\%: ¿son $x_4$ y $x_5$ conjuntamente relevantes? ($F_{0.95,(2,114)}=3.08$)


\textbf{Solución:}
\[
H_0: \beta_4 = \beta_5 = 0 \qquad H_1: \text{al menos uno} \neq 0
\]

Con $q=2$ y $n-k-1 = 120-5-1 = 114$:
\[
F = \frac{(R^2_{UR}-R^2_R)/q}{(1-R^2_{UR})/(n-k-1)} = \frac{(0.78-0.74)/2}{(1-0.78)/114} = \frac{0.02}{0.001930} \approx 10.36
\]

Como $F = 10.36 > 3.08$, \textbf{se rechaza $H_0$}.

\begin{conclusion}
Las variables $x_4$ y $x_5$ son conjuntamente relevantes al 5\%.
\end{conclusion}

\begin{nota}
El dato del $\bar{R}^2$ es un distractor. El test F sobre restricciones múltiples usa $R^2$ (no ajustado). El $\bar{R}^2$ sirve para comparar modelos en términos de ajuste penalizado, pero no reemplaza al test F.
\end{nota}


\item Con $n=80$ observaciones se estima:
\[
\ln(y) = \beta_0 + \beta_1 x_1 + \beta_2 x_2 + \beta_3 x_3 + \beta_4 x_4 + u
\]
Datos: $SST = 500$, $SSR_{UR} = 150$. Al imponer $H_0: \beta_2 = \beta_3 = \beta_4 = 0$ se obtiene $SSR_R = 240$.
\begin{enumerate}
    \item[(a)] Pruebe $H_0$ al 5\%. ($F_{0.95,(3,75)}=2.73$)
    \item[(b)] Calcule el F de significancia global del modelo no restringido. ($F_{0.95,(4,75)}=2.49$)
\end{enumerate}

\textbf{Solución:}

\textbf{(a)} Con SSR:
\[
F = \frac{(SSR_R - SSR_{UR})/q}{SSR_{UR}/(n-k-1)} = \frac{(240-150)/3}{150/75} = \frac{30}{2} = 15
\]
Como $15 > 2.73$, \textbf{se rechaza $H_0$}: las tres variables son conjuntamente significativas.

\textbf{(b)} Para significancia global necesitamos $R^2_{UR}$:
\[
R^2_{UR} = 1 - \frac{SSR_{UR}}{SST} = 1 - \frac{150}{500} = 0.70
\]
\[
F = \frac{R^2/k}{(1-R^2)/(n-k-1)} = \frac{0.70/4}{0.30/75} = \frac{0.175}{0.004} = 43.75
\]
Como $43.75 > 2.49$, \textbf{se rechaza $H_0$}.

\begin{conclusion}
El modelo es globalmente significativo; al menos uno de los regresores explica significativamente $\ln(y)$.
\end{conclusion}

\item Se estima (interpretación aproximada):
\[
\ln(salario) = \beta_0 + \beta_1 exper + \beta_2 urbano + \beta_3 (urbano \times exper) + \beta_4 sindicato + u
\]
Con resultados:
\[
\hat{\beta}_0 = 2.40, \; \hat{\beta}_1 = 0.040, \; \hat{\beta}_2 = 0.12, \; \hat{\beta}_3 = 0.005, \; \hat{\beta}_4 = 0.15
\]
donde $urbano=1$ si reside en zona urbana y $sindicato=1$ si está sindicalizado.
\begin{enumerate}
    \item[(a)] Efecto aproximado de un año adicional de experiencia para un trabajador \emph{rural}.
    \item[(b)] Efecto para un trabajador \emph{urbano}.
    \item[(c)] Interprete $\hat{\beta}_3$.
    \item[(d)] Interprete $\hat{\beta}_4$.
    \item[(e)] Diferencia salarial aproximada entre urbano y rural (no sindicalizados) con $exper=10$.
    \item[(f)] Ecuaciones de predicción para los cuatro subgrupos.
\end{enumerate}
\textbf{Solución:}

\textbf{(a)} Rural ($urbano=0$):
\[
\frac{\partial \ln(salario)}{\partial exper} = \hat{\beta}_1 = 0.040 \Rightarrow \text{aprox. } 4.0\%
\]

\textbf{(b)} Urbano ($urbano=1$):
\[
\frac{\partial \ln(salario)}{\partial exper} = \hat{\beta}_1 + \hat{\beta}_3 = 0.045 \Rightarrow \text{aprox. } 4.5\%
\]

\textbf{(c)} $\hat{\beta}_3 = 0.005$: el retorno de un año adicional de experiencia es aproximadamente 0.5 puntos porcentuales mayor en zona urbana que rural.

\textbf{(d)} $\hat{\beta}_4 = 0.15$: un trabajador sindicalizado gana aproximadamente $15\%$ más que uno no sindicalizado, ceteris paribus.

\textbf{(e)} Diferencial urbano vs rural con $exper=10$, $sindicato=0$:
\[
\Delta \ln(salario) = \hat{\beta}_2 + \hat{\beta}_3 \cdot exper = 0.12 + 0.005 \cdot 10 = 0.17
\]
Un urbano con 10 años de experiencia gana aproximadamente $17\%$ más que un rural comparable.

\textbf{(f)} Ecuaciones por subgrupo:
\begin{itemize}[leftmargin=2em]
    \item \textbf{Rural, no sindicalizado:} $\widehat{\ln(salario)} = 2.40 + 0.040 \cdot exper$
    \item \textbf{Rural, sindicalizado:} $\widehat{\ln(salario)} = 2.55 + 0.040 \cdot exper$
    \item \textbf{Urbano, no sindicalizado:} $\widehat{\ln(salario)} = 2.52 + 0.045 \cdot exper$
    \item \textbf{Urbano, sindicalizado:} $\widehat{\ln(salario)} = 2.67 + 0.045 \cdot exper$
\end{itemize}

\begin{center}
\begin{tabular}{lcc}
\toprule
\textbf{Subgrupo} & \textbf{Intercepto} & \textbf{Pendiente en exper} \\
\midrule
Rural, no sindicalizado & 2.40 & 0.040 ($4.0\%$) \\
Rural, sindicalizado    & 2.55 & 0.040 ($4.0\%$) \\
Urbano, no sindicalizado & 2.52 & 0.045 ($4.5\%$) \\
Urbano, sindicalizado   & 2.67 & 0.045 ($4.5\%$) \\
\bottomrule
\end{tabular}
\end{center}

\item Una cadena de retail estima:
\[
Ventas = \beta_0 + \beta_1 PRECIO + \beta_2 PUBLICIDAD + \beta_3 TAMANO + \beta_4 COMPETENCIA + u
\]
con $n=180$ tiendas repartidas en 4 zonas.
\begin{center}
\begin{tabular}{|c|c|}
\hline
Modelo & SSR \\
\hline
Base (pooled) & 1200 \\
Zona 1 & 180 \\
Zona 2 & 210 \\
Zona 3 & 195 \\
Zona 4 & 215 \\
\hline
\end{tabular}
\end{center}
Teste al 5\% si las ventas dependen de la zona. ($F_{0.95,(15,160)} = 1.73$)


Con $G=4$, $k=4$ regresores y por tanto $k+1=5$ parámetros por modelo:
\[
SSR_{UR} = 180+210+195+215 = 800, \quad SSR_R = 1200
\]
Grados de libertad:
\[
df_1 = (G-1)(k+1) = 3\cdot 5 = 15, \quad df_2 = n - G(k+1) = 180-20 = 160
\]
\[
F = \frac{(SSR_R - SSR_{UR})/15}{SSR_{UR}/160} = \frac{400/15}{800/160} = \frac{26.67}{5} = 5.33
\]
Como $5.33 > 1.73$, \textbf{se rechaza $H_0$}.

\begin{conclusion}
Los parámetros difieren entre las cuatro zonas; un modelo pooled no es adecuado.
\end{conclusion}




\item Con $n=200$ observaciones se estima inicialmente:
\[
y = \beta_0 + \beta_1 x_1 + \beta_2 x_2 + u
\]
Se sospecha que los coeficientes son distintos entre hombres ($D=0$) y mujeres ($D=1$). Para responder esto, utilice el \textbf{enfoque de dummies interactivas} estimando:
\[
y = \alpha_0 + \alpha_1 x_1 + \alpha_2 x_2 + \gamma_0 D + \gamma_1 (D \cdot x_1) + \gamma_2 (D \cdot x_2) + v
\]
Se obtienen $R^2_{UR} = 0.68$ y $R^2_R = 0.62$ (modelo sin los términos con $D$).
\begin{enumerate}
    \item[(a)] Plantee formalmente la hipótesis nula del test.
    \item[(b)] Interprete qué representa cada uno de los coeficientes $\gamma_0$, $\gamma_1$ y $\gamma_2$.
    \item[(c)] Realice el test al 5\%. ($F_{0.95,(3,194)}=2.65$)
    \item[(d)] ¿Qué ventaja tiene el enfoque con dummies interactivas sobre estimar dos regresiones separadas?
\end{enumerate}


\textbf{(a)} Hipótesis:
\[
H_0: \gamma_0 = \gamma_1 = \gamma_2 = 0 \qquad H_1: \text{al menos uno} \neq 0
\]
Bajo $H_0$ no hay diferencias entre grupos (ni en intercepto ni en pendientes), lo que equivale al modelo pooled original.

\textbf{(b)} Interpretación de los coeficientes:
\begin{itemize}[leftmargin=2em]
    \item $\gamma_0$: diferencia en el \textbf{intercepto} entre mujeres y hombres (shift de nivel).
    \item $\gamma_1$: diferencia en la \textbf{pendiente de $x_1$} entre mujeres y hombres.
    \item $\gamma_2$: diferencia en la \textbf{pendiente de $x_2$} entre mujeres y hombres.
\end{itemize}
Así, para hombres ($D=0$): $y = \alpha_0 + \alpha_1 x_1 + \alpha_2 x_2$; y para mujeres ($D=1$):
\[
y = (\alpha_0 + \gamma_0) + (\alpha_1 + \gamma_1) x_1 + (\alpha_2 + \gamma_2) x_2
\]

\textbf{(c)} Con $q=3$ y $n-k-1 = 200-5-1 = 194$:
\[
F = \frac{(R^2_{UR}-R^2_R)/q}{(1-R^2_{UR})/(n-k-1)} = \frac{(0.68-0.62)/3}{(1-0.68)/194} = \frac{0.02}{0.001649} \approx 12.13
\]
Como $12.13 > 2.65$, \textbf{se rechaza $H_0$}.

\begin{conclusion}
Los coeficientes difieren significativamente entre hombres y mujeres al 5\%.
\end{conclusion}

\textbf{(d) Ventajas del enfoque con dummies:}
\begin{enumerate}[leftmargin=2em]
    \item \textbf{Tests parciales:} permite testear subconjuntos específicos de coeficientes (por ejemplo, solo las pendientes manteniendo intercepto común, o solo el intercepto). Con el enfoque de regresiones separadas esto no se puede: o todo varía o nada varía.
    \item \textbf{Errores estándar directos:} los coeficientes $\gamma_j$ vienen con sus propios errores estándar, permitiendo tests $t$ individuales sobre cada diferencia específica.
    \item \textbf{Eficiencia:} al estimar en una sola muestra, se aprovechan todos los datos para estimar los parámetros comunes bajo el supuesto de varianza homogénea.
\end{enumerate}

\end{enumerate}
\end{document}